\documentclass[11pt,a4paper]{article}
\usepackage[utf8]{inputenc}
\usepackage[T1]{fontenc}
\usepackage[margin=2.2cm]{geometry}
\usepackage{amsmath,amssymb,booktabs,array}
\usepackage{enumitem}
\usepackage{hyperref}
\usepackage{xcolor}

\hypersetup{
  colorlinks=true,
  linkcolor=blue!50!black,
  urlcolor=blue!50!black
}

\title{Validated Lab Guide\\Envelope Modulation of Random Signals}
\author{Signal Theory Workshop 02}
\date{April 2026}

\begin{document}
\maketitle

\section*{Purpose}
This guide is the validated and cleaned deliverable requested by the repository
instructions. The refined brief in \texttt{docs/lab00\_.pdf} is treated as the
main source of truth. The result preserves the exact stage names and the same
subscript notation used in the original PDF while tightening the structure
around parameters, estimators, and required outputs.

\section*{Objective}
Study how a white Gaussian input process evolves through a conventional AM
communication chain:
\begin{enumerate}[leftmargin=1.5em]
  \item baseband generation by RC low-pass filtering,
  \item conventional AM modulation,
  \item envelope detection by rectification plus low-pass filtering.
\end{enumerate}

At every stage, estimate the autocorrelation function (ACF) and power spectral
density (PSD), then compare the recovered signal with the original message.

\section*{Signal Model}
Let $x[n]$ be discrete-time white Gaussian noise:
\[
  x[n] \sim \mathcal{N}(0,1).
\]

The filtered message signal is denoted by $x_1[n]$. The conventional AM signal
is
\[
  x_2[n] = (A_0 + x_1[n]) \cos(2\pi f_0 nT_s),
\]
with $T_s = 1/f_s$. The peak and RMS modulation indices are
\[
  \mu_{peak} = \frac{\max_n |x_1[n]|}{A_0},
  \qquad
  \mu_{rms} = \frac{\sigma_{x_1}}{A_0}.
\]
To avoid overmodulation on the observed record, require
\[
  \mu_{peak} < 1
  \qquad \Longleftrightarrow \qquad
  \min_n (A_0 + x_1[n]) > 0.
\]

Under the cycle-averaged interpretation from the lab statement, the AM ACF is
\[
  R_{x_2x_2}(\tau)
  \approx
  \left(
    \frac{A_0^2}{2}
    +
    \frac{R_{x_1x_1}(\tau)}{2}
  \right)
  \cos(2\pi f_0 \tau),
\]
and the AM PSD contains the carrier and translated sidebands around $\pm f_0$.

\section*{Parameter Selection}
Use the recommended values from the original PDF, with validated adjustments:
the executed notebook uses $A_0 = 6.0$ instead of $A_0 = 4.0$ because the
original suggestion can violate $\mu_{peak} < 1$ on a long normalized Gaussian
record, and it uses $f_{demod} = 300\,\text{Hz}$ so the discrete-time full-wave
detector preserves the recovered envelope while still attenuating the main
rectifier ripple around $2f_0 = 600\,\text{Hz}$.
\begin{center}
\begin{tabular}{>{\raggedright\arraybackslash}p{5.3cm} p{6.5cm}}
\toprule
Baseband cutoff & $f_n = 10\,\text{Hz}$ \\
Carrier ratio & $k = 30$ \\
Carrier frequency & $f_0 = k f_n = 300\,\text{Hz}$ \\
Carrier amplitude & $A_0 = 6.0$ \\
Sampling frequency & $f_s = 15{,}000\,\text{Hz}$ \\
Observation time & $T = 50\,\text{s}$ \\
Number of samples & $N = f_s T = 750{,}000$ \\
Demodulator cutoff & $f_{demod} = 300\,\text{Hz}$ \\
PSD method & Welch, Hann window, $N_{seg}=16384$, 50\% overlap \\
\bottomrule
\end{tabular}
\end{center}

\section*{Required Stages}
The Colab notebook must preserve these stage headings exactly:

\subsection*{Stage 0: Generation of Input Signal}
\begin{itemize}[leftmargin=1.5em]
  \item Generate $x[n] \sim \mathcal{N}(0,1)$.
  \item Estimate and plot $R_{xx}[\ell]$.
  \item Estimate and plot $S_{xx}(f)$.
\end{itemize}

\subsection*{Stage 1: RC Low-Pass Filtering}
\begin{itemize}[leftmargin=1.5em]
  \item Filter $x[n]$ with a first-order RC-equivalent low-pass filter.
  \item Obtain the message signal $x_1[n]$ and normalize it to unit variance.
  \item Estimate and plot $R_{x_1x_1}[\ell]$ and $S_{x_1x_1}(f)$.
  \item Compare against the reference RC shapes
  \[
    R_{x_1x_1}(\tau) \propto e^{-2\pi f_n |\tau|},
    \qquad
    S_{x_1x_1}(f) \propto \frac{1}{1 + (f/f_n)^2}.
  \]
\end{itemize}

\subsection*{Stage 2: Conventional AM Modulation}
\begin{itemize}[leftmargin=1.5em]
  \item Form
  \[
    x_2[n] = (A_0 + x_1[n]) \cos(2\pi f_0 nT_s).
  \]
  \item Verify $\mu_{peak} < 1$ and report $\mu_{rms}$.
  \item Plot the AM waveform with its envelope.
  \item Estimate and plot $R_{x_2x_2}[\ell]$ and $S_{x_2x_2}(f)$.
\end{itemize}

\subsection*{Stage 3: Envelope Detection}
\begin{itemize}[leftmargin=1.5em]
  \item Rectify: $x_{rect}[n] = |x_2[n]|$.
  \item Low-pass filter the rectified waveform to obtain $x_{env}[n]$.
  \item Remove the DC component and normalize the result to obtain $x_{demod}[n]$.
  \item Estimate and plot $R_{demod}[\ell]$ and $S_{demod}(f)$.
  \item Compute
  \[
    \rho = \mathrm{corr}(x_1, x_{demod}),
    \qquad
    \mathrm{NMSE} =
    \frac{\mathbb{E}\!\left[(x_1 - x_{demod})^2\right]}{\sigma_{x_1}^2},
    \qquad
    \mathrm{SNR} = -10 \log_{10}(\mathrm{NMSE})\ \mathrm{dB}.
  \]
\end{itemize}

\subsection*{Comparison and Analysis}
\begin{itemize}[leftmargin=1.5em]
  \item Compare $x_1[n]$ and $x_{demod}[n]$ in time and frequency.
  \item Compare $R_{x_1x_1}[\ell]$ with $R_{demod}[\ell]$.
  \item Compare $S_{x_1x_1}(f)$ with $S_{demod}(f)$.
  \item Discuss sidebands, envelope positivity, and demodulation error.
\end{itemize}

\section*{Required Deliverables}
The submission must contain:
\begin{enumerate}[leftmargin=1.5em]
  \item the executed notebook,
  \item this PDF guide,
  \item the following figures inside the notebook:
  \begin{itemize}[leftmargin=1.5em]
    \item Stage 0 ACF and PSD,
    \item Stage 1 ACF and PSD with theoretical overlays,
    \item Stage 2 waveform, ACF, and PSD,
    \item Stage 3 detection and recovery plots,
    \item modulation-index analysis with four subplots,
    \item error analysis with four subplots,
    \item comprehensive comparison with a $2\times2$ layout.
  \end{itemize}
  \item the reported metrics:
  \begin{itemize}[leftmargin=1.5em]
    \item $\mu_{peak}$ and $\mu_{rms}$,
    \item $\rho$,
    \item NMSE,
    \item SNR in dB,
    \item power efficiency
    \[
      \eta = \frac{\mu_{rms}^2}{1 + \mu_{rms}^2} \times 100\%.
    \]
  \end{itemize}
\end{enumerate}

\section*{Validation Notes}
This guide intentionally keeps the original notation and stage names from the
lab PDF while making the contract explicit:
\begin{itemize}[leftmargin=1.5em]
  \item the AM chain is conventional AM, not DSB-SC,
  \item the detector is full-wave rectification plus low-pass filtering,
  \item the comparison signal is $x_1[n]$,
  \item the recovered output is $x_{demod}[n]$ after DC removal and normalization.
\end{itemize}

\end{document}
